\documentclass[11pt,a4]{article}
\usepackage[francais]{babel}
\usepackage[utf8]{inputenc}
\usepackage{graphics}

\textheight=17cm
\sloppy
\usepackage{wallpaper}
\usepackage{hyperref}
\usepackage[margin=1in,footskip=0.25in]{geometry}
\usepackage{algpseudocode}
\usepackage{tikz}

\begin{document}
\CenterWallPaper{1}{Modelecourrier2012.pdf}

%Bonjour les amis.
\medskip
\medskip
\vspace*{0.35cm}
\noindent\textbf{Internship subject}: \emph{Cache efficient binary search}\newline
%\textbf{Advisors}: \href{https://surakuma.github.io}{Suraj kumar} and \href{https://perso.ens-lyon.fr/bora.ucar}{Bora Uçar}\newline
\textbf{Internship duration}: 5-6 months \newline
\textbf{Location}: ROMA team, LIP, ENS Lyon, France

\section*{Objectives}
\begin{itemize}
	\item[$\bullet$] Study bounds on the minimum number of cache misses on a sequential machine for binary search
%	 and design approaches that achieve the obtained bounds
	\item[$\bullet$] Design an approach to store elements of an array that improves the cache utilization for binary search on a sequential machine
\end{itemize}
\section*{Context}
Each element of an array may not be accessed with the same probability in binary search. Let us consider the following array $A$ with $n=15$ elements, 1 through 15.
\[
\begin{array}{|*{15}{c|}}
\hline
1 & 2 & 3 & 4 & 5 & 6 & 7 & 8 & 9 & 10 & 11 & 12 & 13 & 14 & 15 \\
\hline
\end{array}
\] 
\begin{algorithmic}
	\Require $A$ is an array with $n$ elements and the search element is $key$.
	\Ensure If $key$ is found in $A$ then its position, otherwise $-1$.
\end{algorithmic}
\begin{minipage}{0.4\linewidth}
\begin{algorithmic}
%	\Require $A$ is an array with $n$ elements and the search element is $key$.
%	\Ensure If $key$ is found in $A$ then its position, otherwise $-1$.
	\State start = $0$
	\State end = $n-1$
	\While{start $\leq$ end}
	\State mid = (start+end)/2
	\If{A[mid]== key}
	\State key is found, \Return mid
	\ElsIf{A[mid]$\ge$ key}
	\State end = mid-1
	\Else
	\State start = mid +1
	\EndIf
	\EndWhile
	\State \Return -1
\end{algorithmic}
\end{minipage}
$\quad$
\begin{minipage}{0.56\linewidth}
\begin{tikzpicture}[
scale=0.625,
level distance=1.2cm,
level 1/.style={sibling distance=7cm},
level 2/.style={sibling distance=3.5cm},
level 3/.style={sibling distance=1.7cm},
every node/.style={circle, draw, minimum size=7mm, transform shape}
]
\node {8}
child {node {4}
	child {node {2}
		child {node {1}}
		child {node {3}}
	}
	child {node {6}
		child {node {5}}
		child {node {7}}
	}
}
child {node {12}
	child {node {10}
		child {node {9}}
		child {node {11}}
	}
	child {node {14}
		child {node {13}}
		child {node {15}}
	}
};
\end{tikzpicture}
\end{minipage}

\medskip

The element $8$ is always accessed irrespective of the $key$ element. Probability of accessing the element $4$ is $probability(key\ne 8)\cdot\frac{1}{2}$. The element $12$ is also accessed with the same probability. Similarly, the elements $2,6,10,14$ have same probability of getting accessed. The elements $1,3,5,7,9,11,13,15$ are accessed with minimum probability.

\newpage

\section*{Approach to improve cache utilization}
When $i$th element is accessed and it is not equal to the $key$ then the sum of probabilities of accessing the elements $i-(k/4+1)$ and $i+(k/4+1)$ is $1$. Here $k$ is the number of active (non-discarded) elements in the current search space. For example, when the element $12$ is accessed and it is not equal to the key, then $k=6$ and the sum of probabilities to access $10$ and $14$ is $1$. We can make better utilization of cache by storing the elements $i$, $i-(k/4+1)$ and $i+(k/4+1)$ together.  
\section*{Main Activities}

We will perform the following activities:
\begin{itemize}
	\item[$\bullet$] Determine the bound on the minimum number of cache misses in the worst case for binary search
	\item[$\bullet$] Design an approach and/or a data structure that achieves the lower bound for all inputs
\end{itemize}

\section*{Example}
Let us assume that the cache is empty in the beginning and each cache line can contain $3$ elements. 

There are 5 cache lines for our example. If we store elements contiguously, then there are $3$ cache misses in the worst case scenario. If the $key=15$, then $3$ cache lines will be accessed. Hence the total number of cache misses is $3$. 

\[
\begin{array}{|c|c|c|}
\hline
1 & 2 & 3 \\ \hline
4 & 5 & 6 \\ \hline
7 & 8 & 9 \\ \hline
10 & 11 & 12 \\ \hline
13 & 14 & 15 \\ \hline
\end{array}
\]




Now consider the following cache lines. The number of cache misses for any $key$ is at max $2$.  
\[
\begin{array}{|c|c|c|}
\hline
8 & 4 & 12 \\ \hline
2 & 1 & 3 \\ \hline
6 & 5 & 7 \\ \hline
10 & 9 & 11 \\ \hline
14 & 15 & 16 \\ \hline
\end{array}
\]







\section*{Skills}
Familiarity with tree data structure will be much appreciated. 

\bibliographystyle{plain}
%\footnotesize \bibliography{mttkrp-gpu}




\end{document}